\documentclass[11pt,letterpaper]{article}
\usepackage[T1]{fontenc}
\usepackage{lmodern}
\usepackage[margin=1in]{geometry}
\usepackage{amsmath,amssymb,amsthm,mathtools,microtype}
\usepackage{booktabs,array,enumitem,xurl,xcolor}
\usepackage[colorlinks=true,linkcolor=blue!50!black,citecolor=blue!50!black,urlcolor=blue!50!black]{hyperref}
\hypersetup{pdftitle={Finite-information Lyapunov potentials: the depth-one infimum},pdfauthor={Jake S. Brusca}}
\pdfinfoomitdate=1
\pdftrailerid{}
\pdfsuppressptexinfo=15
\setlength{\emergencystretch}{1.5em}
\setlength{\parskip}{0.25em}
\newtheorem{theorem}{Theorem}
\newtheorem{lemma}{Lemma}
\newcommand{\ind}{\mathbf{1}}
\newcommand{\NN}{\mathbb{N}}
\newcommand{\RR}{\mathbb{R}}
\title{Finite-information Lyapunov potentials for the Syracuse map\\[0.3em]\large The exact depth-one infimum: a written research draft}
\author{Jake S. Brusca}
\date{27 November 2026}
\begin{document}
\maketitle
\begin{abstract}
For the odd-only Syracuse map, consider potentials of the form
$\log x+f(v_2(x+1))$, with $f$ an arbitrary real-valued function on positive
integers. We give a complete written argument for existence of each potential's
relative odd-seed failure density, an exact series for its infimum, and
nonattainment. The argument partitions initial seeds into valuation cylinders;
it assumes neither Collatz convergence nor independence along an integer orbit.
An exact rational enclosure has width $2^{-81}$. This addresses only the
depth-one subquestion of Polson--Sokolov O3$'$, not its arbitrary-depth part.
The manuscript remains a research draft with assurance limits stated below.
\end{abstract}

\begin{quote}\small
\textbf{Personal project; work in progress.} I am not an expert in this field.
I am exploring these ideas for fun and to learn. Corrections, counterexamples
and suggestions are welcome. ChatGPT assisted with ideas, proofs, code,
writing and review. Separate AI contexts examined the source model and proof;
this is not human expert refereeing, peer acceptance or a priority claim.
Nothing here proves ordinary Collatz convergence or global forced capture.
\end{quote}

\paragraph{Source and formal-assurance snapshot.}
The source problem is O3$'$ in Polson and Sokolov~\cite{polson}, edition~1,
version~1.4.0, label \texttt{prob:ch11:o3p}. Its exact source commit is
\href{https://github.com/VadimSokolov/collatz/commit/5840497a23fc0d64127b850aea68862d7c93f2b7}{\nolinkurl{5840497a23fc0d64127b850aea68862d7c93f2b7}}.
The block convention is in the same version's
\href{https://github.com/VadimSokolov/collatz/blob/5840497a23fc0d64127b850aea68862d7c93f2b7/book/chapters/ch05c_mmcollatz.tex}{block-majorization chapter}.
The public \path{lean/external_analysis/} package contains proofs of all
fourteen frozen depth-one targets, using the Lean mathematical library
\cite{mathlib}. The historical mathlib paper describes the library; the
release lockfile pins the actual contemporary dependency revision.
At source commit \texttt{8524770}, the dedicated analytic CI job passed a
clean build, exhaustive axiom audit, source-bound sandbox replay, frozen-target
comparison, all six bundled checker modes and all nine hostile fixtures.
Exact source hashes and receipts are in \path{verification/external_analysis/}.
Later manuscript changes still require final-tree regression checks. Ordinary compilation
is not independent kernel replay or human review. Consult the release's
per-declaration coverage register; this paper does not claim those pending
assurance gates have passed.

\section{Related work and scope of attribution}
Finite parity data for Collatz iteration have classical residue-class
realizations. Terras~\cite{terras}, Section~1, develops finite parity-vector
coding and its counting consequences. Bernstein and Lagarias~\cite{bernstein-lagarias}
study the conjugacy between the half-step Collatz map and the binary shift
on the $2$-adic integers. The accelerated odd map here is obtained by grouping
steps; its joint exit/block cylinders and relative odd-seed normalization are
proved explicitly below. These population laws do not make the successive
symbols of one fixed integer trajectory independent.

The potential-cycle argument is related to classical finite graph
optimization: Karp~\cite{karp} characterizes the minimum cycle mean of a
finite digraph. Adding a difference of vertex potentials leaves every cycle
sum unchanged. Here the graph of run lengths is countable, the objective
weights failed edges rather than minimizing mean cycle weight, and the
correction may be unbounded. The infinite infimum, density transfer and
nonattainment therefore require the separate arguments in this paper.

The Diophantine step uses Dirichlet approximation~\cite{dirichlet}.
A two-sided small fractional error alone is insufficient for our positive
return indices: the nonattainment proof below explicitly converts a negative
error into a positive near-integer return. The exact formula here is a
candidate contribution to the depth-one question of Polson and Sokolov,
not a claim that parity coding, cycle obstructions or irrational approximation
are new. The references are a scoped comparison, not a completed novelty
review.

\section{The precise depth-one question}
On positive odd integers define
\[
 K(x)=v_2(3x+1),\qquad S(x)=\frac{3x+1}{2^{K(x)}},\qquad r(x)=v_2(x+1).
\]
A block consists of the consecutive strides with $K=1$ followed by the first
stride with $K\ge2$. Its remaining length at $x$ is $r(x)$, including that
last stride. A depth-one correction reads only this remaining length:
\[
 f:\NN_{\ge1}\longrightarrow\RR,\qquad
 V_f(x)=\log x+f(r(x)).
\]
There is no boundedness, support or regularity condition on $f$. All logarithms
in the potential are natural. Failure means nondecrease, including equality:
\[
 V_f(S(x))\ge V_f(x).
\]
Write $d(f)$ for the limit, when it exists, of the number of failures at
positive odd $x\le X$ divided by the number of positive odd $x\le X$.
Equivalently, count $x=2n+1$ for $0\le n<N$ and divide by $N$.

\begin{theorem}[Depth-one claim proved in the written argument]
For every such $f$, the failure density $d(f)$ exists. Set
$a=\log(3/2)$, $b=\log2$, and
\[
 K_r=\left\lfloor1+\frac{ra}{b}\right\rfloor,\qquad
 T_r=2^{-r}\bigl(\tfrac12-2^{-K_r}\bigr).
\]
Then
\begin{equation}\label{eq:infimum}
 p_1:=\inf_f d(f)=\sum_{r\ge1}T_r.
\end{equation}
No correction attains $p_1$. For every integer $M\ge0$,
\begin{equation}\label{eq:tail}
 0\le p_1-\sum_{r=1}^{M}T_r\le2^{-M-1}.
\end{equation}
\end{theorem}

This does not determine $p_h$ for every depth $h$, prove an inverse-depth
rate, or give higher-depth closed forms. These are separate parts of O3$'$.

\section{An exact partition of initial populations}
If $r(x)=r\ge2$, write $x=2^ru-1$ with $u$ positive odd. Then
\[
 3x+1=2(3\cdot2^{r-1}u-1),\qquad
 K(x)=1,\qquad S(x)+1=3\cdot2^{r-1}u.
\]
Consequently $r(S(x))=r-1$. This ascending edge $r\to r-1$ is one odd
residue class modulo $2^{r+1}$ and has relative odd density $2^{-r}$.
In the parameter $x=2n+1$, its exact class is
\begin{equation}\label{eq:ascending-class}
 n\equiv2^{r-1}-1\pmod{2^r}.
\end{equation}

If $r(x)=1$, then $K(x)\ge2$. Fix an exit valuation $k\ge2$ and next length
$r\ge1$. Let $M=2^{k+r}$ and $C=2^{k-1}(2^r-1)$. The joint event is exactly
\begin{equation}\label{eq:exit-class}
 K(2n+1)=k,\quad r(S(2n+1))=r
 \quad\Longleftrightarrow\quad 3n+2\equiv C\pmod M.
\end{equation}
Indeed $3x+1=2(3n+2)=2^k y$, with $y=S(x)$, and the required next length is
$y\equiv2^r-1\pmod{2^{r+1}}$. Conversely the congruence gives
$3n+2=C+qM$ for $q\ge0$, so
\[
 y=2^r-1+q2^{r+1}
\]
is positive odd, has $v_2(y+1)=r$, and satisfies $3x+1=2^ky$. Oddness makes
the valuation exactly $k$. Multiplication by $3$ is invertible modulo $M$;
hence this event is one class in $n$ modulo $2^{k+r}$ and has relative odd
density $2^{-k-r}$. This is not a conditional probability given an exit.

For completeness, a class $n\equiv c\pmod M$, $0\le c<M$, has exactly
\[
 \left\lfloor\frac{N+M-1-c}{M}\right\rfloor
\]
representatives in $0\le n<N$. Its count differs from $N/M$ by at most one,
which proves the asserted densities directly. Total-integer density would
be half the relative odd density and is not the normalization used here.

All source seeds are partitioned into ascending edges $r\to r-1$,
$r\ge2$, and exits $1\to r$ with $k\ge2$. Their weights sum to one:
\[
 \sum_{r\ge2}2^{-r}+\sum_{r\ge1,\ k\ge2}2^{-r-k}
 =\tfrac12+\tfrac12=1.
\]
This is a statement about initial seed populations. It assumes no trajectory
convergence. The fixed point $x=1$ is the exit $(k,r)=(2,1)$ and always fails,
but the singleton $\{1\}$ has density zero.

\section{Symbolic gains and existence of every density}
The exact increment is
\[
 \log\frac{S(x)}x=\log3-K(x)\log2+
 \underbrace{\log\!\left(1+\frac1{3x}\right)}_{\varepsilon(x)>0}.
\]
Here $\varepsilon(x)\to0$ as $x\to\infty$. Define the constant symbolic gains
on the partition's edges by
\[
 A_r=a+f(r-1)-f(r),\qquad
 E_{r,k}=a-(k-1)b+f(r)-f(1).
\]
For each fixed cylinder, the actual failure indicator eventually equals
$\ind_{\{A_r\ge0\}}$ or $\ind_{\{E_{r,k}\ge0\}}$.
If the symbolic gain is zero, the strictly positive $\varepsilon(x)$ means
every seed in that cylinder fails. Thus replacing $\ge$ by $>$ would change
the intended symbolic boundary.

Fix $\epsilon>0$. Choose finitely many edge cylinders whose total density
exceeds $1-\epsilon$. On each of these cylinders the failure indicator
stabilizes outside finitely many seeds. Their finite union therefore has its
predicted selected density. Its complement is the complement of a finite
dyadic union and has density below $\epsilon$. Both the limiting lower and
upper densities of the entire failure set lie between the finite selected
weight and that weight plus $\epsilon$. Letting $\epsilon\downarrow0$ proves
existence and the exact formula
\begin{equation}\label{eq:density}
 d(f)=\sum_{r\ge2}2^{-r}\ind_{\{A_r\ge0\}}
 +\sum_{r\ge1,\ k\ge2}2^{-r-k}\ind_{\{E_{r,k}\ge0\}}.
\end{equation}
No uniform stabilization threshold across all edges is required: first retain
finitely many edges, then take the maximum of their thresholds. This order of
quantifiers permits arbitrary unbounded $f$.

\section{A universal lower bound from telescoping cycles}
Follow an exit $1\to r$ with valuation $k$, then the ascending edges back to
$1$. The correction terms telescope:
\begin{equation}\label{eq:cycle}
 E_{r,k}+\sum_{j=2}^{r}A_j=ra-(k-1)b.
\end{equation}
This cycle gain is never zero: equality would give
$3^r=2^{r+k-1}$, contradicting prime factorization. It is positive precisely
when $2\le k\le K_r$. If all ascending edges through $r$ are strictly good,
every such exit must fail. The total mandatory exit weight at height $r$ is
\[
 \sum_{k=2}^{K_r}2^{-r-k}
 =2^{-r}\bigl(\tfrac12-2^{-K_r}\bigr)=T_r.
\]
When $K_r=1$, this is the empty sum and both sides are zero.

If an ascending edge fails, let $L\ge2$ be its least source. Every positive
cycle of height $r<L$ then forces its own exit to fail, while the ascending
edge $L$ fails. These disjoint events give
\[
 d(f)\ge q(L):=2^{-L}+\sum_{r<L}T_r.
\]
If no ascending edge fails, every mandatory exit fails, so
$d(f)\ge q_\infty:=\sum_{r\ge1}T_r$.
Direct subtraction gives
\[
 q(L)-q(L+1)=2^{-L-K_L}>0.
\]
Moreover $q(L)\to q_\infty$, since $2^{-L}\to0$ and the series converges
by $0\le T_r\le2^{-r-1}$. Thus every finite $L$ has
$q(L)>q_\infty$, proving the universal bound $d(f)\ge q_\infty$ and a strict
bound whenever an ascending edge fails.

\section{Approaching the bound and controlling the tail}
For $c>a$, take $f_c(r)=c(r-1)$. All ascending gains are $a-c<0$. The exit gain is
\[
 E_{r,k}=ra-(k-1)b+(r-1)(c-a).
\]
As $c\downarrow a$, every exit indicator tends to the mandatory indicator:
a positive or zero limiting gain stays a failure, and a negative gain is
eventually negative. In the present source model the cycle gain is nonzero,
but this convergence argument does not need that fact. Exit weights are
summable, with total $1/2$, so finite truncation followed by a summable tail
bound gives $d(f_c)\to q_\infty$. Together with the lower bound this proves
\eqref{eq:infimum}.

Each $K_r\ge1$, so $0\le T_r<2^{-r-1}$. Consequently
\[
 0\le\sum_{r>M}T_r\le\sum_{r>M}2^{-r-1}=2^{-M-1},
\]
which proves \eqref{eq:tail}; the upper tail bound is in fact strict.
Finite thresholds can be checked without transcendental approximations:
$ra>(k-1)b$ is equivalent to $3^r>2^{r+k-1}$.

\section{Why the infimum is not attained}
The case of a failed ascending edge already gives $d(f)>p_1$. Otherwise all
$A_j<0$. Put
\[
 \delta=f(2)-f(1)-a>0.
\]
For every $r\ge2$, summing the strictly positive increment surpluses yields
\begin{equation}\label{eq:surplus}
 f(r)-f(1)-(r-1)a\ge\delta.
\end{equation}
The number $\alpha=a/b$ is irrational, since a rational identity
$\alpha=m/n$ would give $3^n=2^{m+n}$.

We need positive integers $r$, arbitrarily large, with $\{r\alpha\}$
arbitrarily close to one. Pigeonhole approximation gives $q>0$ for which
the distance from $q\alpha$ to an integer is below any prescribed
$\epsilon$. If its fractional part is near one, use $r=q$. Otherwise write
its small positive fractional part as $\gamma$ and take
$m=\lfloor1/\gamma\rfloor$. Then $1-\gamma<m\gamma<1$; irrationality
excludes equality. Hence $\{mq\alpha\}$ is near one. Repeating at smaller
tolerances yields arbitrarily large $r$, because a fixed finite set of
positive $r$ has a strictly positive minimum distance of $\{r\alpha\}$
from one.

Choose $r\ge2$ with $b(1-\{r\alpha\})<\delta$ and set $k=K_r+1$. Then
$k-1=\lceil r\alpha\rceil$ and the cycle gain is negative but satisfies
\[
 ra-(k-1)b=-b(1-\{r\alpha\})>-\delta.
\]
Its exit gain adds the surplus in \eqref{eq:surplus}, so it is strictly
positive. This exit therefore fails in addition to every mandatory exit.
It is disjoint from them and has positive weight $2^{-r-k}$. Formula
\eqref{eq:density} gives $d(f)>p_1$. No arbitrary correction attains the infimum.

\section{The exact numerical enclosure and its scope}
Let
\begin{align*}
 L&=\frac{24353541309990355464286619961554936997}
 {170141183460469231731687303715884105728},\\[0.3em]
 U&=\frac{24353541309990355464286690330299114661}
 {170141183460469231731687303715884105728}.
\end{align*}
The exact first eighty terms sum to $L$. The geometric remainder gives
$L\le p_1<U$, with $U-L=2^{-81}$. Both endpoints round at ordinary machine
precision to approximately $0.14313725116205436$; the rational endpoints,
not that decimal, are the enclosure.

A small exact logarithmic bracket suffices to certify all eighty floors:
\[
 \frac{31}{53}<\alpha<\frac{24}{41}.
\]
The lower inequality follows from $2^{84}<3^{53}$ and the upper from
$3^{41}<2^{65}$, by taking logarithms. For each $1\le r\le80$, the integer
certificate $F_r$ satisfies
\[
 F_r\le1+\frac{31r}{53},\qquad
 1+\frac{24r}{41}\le F_r+1.
\]
Strictness of the logarithmic bracket proves $K_r=F_r$, even at the rational
boundary cases $r=53$ and $r=41$. Exact rational summation and the tail then
certify the enclosure. A corrupted floor entry is rejected by the checker.
The finite sum certificate alone encloses a series; identifying that series
with the infimum requires the preceding universal proof and its source model.

\paragraph{Reproduction and remaining assurance.}
The release's \path{certificates/external_problems/verify.py} replays finite
arithmetic. The accompanying analysis proof register must distinguish actual
literal density existence, the infimum identity and nonattainment from
supporting finite certificates. Ordinary Lean builds, transitive axiom audits,
trusted statement comparisons, fresh replay and independent checker coverage
are different gates; final receipts, rather than theorem counts, govern their
status. None supplies human expert refereeing or acceptance by the literature.

\begin{thebibliography}{9}
\bibitem{polson}
N.~Polson and V.~Sokolov.
\emph{On the Collatz Conjecture: Martingales, Momentum and Continued Fractions}.
Edition~1, version~1.4.0, 2026.
\href{https://doi.org/10.5281/zenodo.24113412}{doi:10.5281/zenodo.24113412}.
Problem O3$'$, label \texttt{prob:ch11:o3p};
\href{https://github.com/VadimSokolov/collatz/blob/5840497a23fc0d64127b850aea68862d7c93f2b7/book/chapters/ch11_open.tex}{pinned original source}.
The authors' \href{https://vsokolov.org/books/collatz/results.html}{results index}
was checked on 27 November 2026. An open listing is not a claim of novelty.
\bibitem{mathlib}
The mathlib Community. The Lean mathematical library.
\emph{Proceedings of CPP 2020} (2020), 367--381.
\href{https://doi.org/10.1145/3372885.3373824}{doi:10.1145/3372885.3373824};
\href{https://arxiv.org/abs/1910.09336}{author manuscript}.
\bibitem{terras}
R.~Terras. A stopping time problem on the positive integers.
\emph{Acta Arithmetica} \textbf{30}(3) (1976), 241--252.
\href{https://doi.org/10.4064/aa-30-3-241-252}{doi:10.4064/aa-30-3-241-252}.

\bibitem{bernstein-lagarias}
D.~J.~Bernstein and J.~C.~Lagarias. The $3x+1$ conjugacy map.
\emph{Canadian Journal of Mathematics} \textbf{48}(6) (1996), 1154--1169.
\href{https://doi.org/10.4153/CJM-1996-060-x}{doi:10.4153/CJM-1996-060-x}.

\bibitem{karp}
R.~M.~Karp. A characterization of the minimum cycle mean in a digraph.
\emph{Discrete Mathematics} \textbf{23}(3) (1978), 309--311.
\href{https://doi.org/10.1016/0012-365X(78)90011-0}{doi:10.1016/0012-365X(78)90011-0}.

\bibitem{dirichlet}
P.~G.~Lejeune Dirichlet.
Verallgemeinerung eines Satzes aus der Lehre von den Kettenbr\"uchen
nebst einigen Anwendungen auf die Theorie der Zahlen (1842).
In \emph{Werke}, vol.~I, chapter~XXXV, pp.~633--638.
\href{https://doi.org/10.1017/CBO9781139237338.037}{doi:10.1017/CBO9781139237338.037}.

\end{thebibliography}
\end{document}
